\usepackage[utf8]{inputenc}
\usepackage[T1]{fontenc}
% \usepackage[slovak]{babel}
\usepackage[english]{babel}
\usepackage[a4paper]{geometry}
\usepackage[
    left = \glqq,% 
    right = \grqq,% 
    leftsub = \glq,% 
    rightsub = \grq%
]{dirtytalk}

% SETTINGS - NAMES
\newcommand{\myTitle}[0] {Deep Learning Methods for Perception and Prediction in Autonomous Systems}
\newcommand{\myTitleSK}[0] {Metódy hlbokého učenia pre percepciu a predikciu v autonómnych systémoch}
\newcommand{\myName}[0] {Bc. Richard Čerňanský}
\newcommand{\mySupervisor}[0] {Ing. Matej Halinkovič, PhD.}
\newcommand{\myEvidenceNumber}[0] {FIIT-XXXX-XXXXX}
\newcommand{\myDate}[0] {May 20xx}
\newcommand{\myDateSK}[0] {Máj 20xx}
\newcommand{\myStudyProgram}[0] {Intelligent Software Systems}
\newcommand{\myStudyProgramSK}[0] {Inteligentné softvérové systémy}
\newcommand{\myStudyField}[0] {Computer Science}
\newcommand{\myStudyFieldSK}[0] {Informatika}
\newcommand{\myUniversity}[0] {Slovak University of Technology in Bratislava}
\newcommand{\myUniversitySK}[0] {Slovenská technická univerzita v Bratislave}
\newcommand{\myFaculty}[0] {Faculty of Informatics and Information Technologies}
\newcommand{\myFacultySK}[0] {Fakulta informatiky a informačných technológií}
\newcommand{\myInstitute}[0] {Institute of Informatics, Information Systems and Software Engineering}
\newcommand{\myInstituteSK}[0] {Ústav informatiky, informačných systémov a softvérového inžinierstva}

\usepackage[parfill]{parskip}
\usepackage{enumitem}

\usepackage{graphicx}
\usepackage{float}
\usepackage{longtable}
\usepackage{setspace}

% Spacing
\setstretch{\mySpacing}

\setcounter{secnumdepth}{3}
\setcounter{tocdepth}{3}

\usepackage{tabularx}
\newsavebox\mybox
\usepackage{fancyhdr}
\pagestyle{fancy}
\lhead{\nouppercase{\leftmark}}
\chead{}
\rhead{}
\lfoot{}
\cfoot{\thepage}
\rfoot{}


%style=iso-numeric
%style=authortitle-dw
\usepackage[backend=biber,sorting=none]{biblatex}
\defbibheading{references}[Literatúra]{ 
  \chapter*{#1}
  \markboth{#1}{#1}
}
\defbibheading{referencessec}[Literatúra]{ 
  \section*{#1}
  \markboth{#1}{#1}
}
\bibliography{\myBibliography}


% Listing as figure
%\usepackage{libs/minted}
%\usepackage[section]{minted}

\usepackage{listing}

% openright does not work :(
\let\tmp\oddsidemargin
\let\oddsidemargin\evensidemargin
\let\evensidemargin\tmp
\reversemarginpar

\usepackage{lscape}
\usepackage{afterpage}

\usepackage{lipsum}

% Acronyms in automatic way
% \usepackage{acro}[=v2] % default is v3, and causing timeout
% \acsetup{list/template=longtable}


% GN setting begin
    % https://tug.org/FontCatalogue/
    % \usepackage{mathpazo}           % EN + math font in one
    \usepackage{newpxtext}        % slovencina, aby nerobila problem s CRZP
    \usepackage{eulerpx}          % math font for newpxtext
    \let\Bbbk\relax              % avoid clash: newpxmath (via eulerpx) already defines \Bbbk
    \usepackage{amsfonts, amssymb, amsmath}     % math symboly
    
    % URL without red/blue boxing
    \usepackage[hidelinks,unicode]{hyperref}
    
    % Clean the space among items
    \usepackage{etoolbox}
    \AtBeginEnvironment{itemize}{\setlength\parskip{0pt}}
    \AtBeginEnvironment{enumerate}{\setlength\parskip{0pt}}
        
    % include PDF zadanie z AISu (alebo draft PDF zadania z YonBanu)
    \usepackage{pdfpages}

    % Acronyms/Abbreviations in automatic way
    \usepackage{longtable}
    \usepackage{acro} % v2 rollback is broken on current LaTeX; use current v3
    % https://www.overleaf.com/learn/latex/Glossaries

%\usepackage[acronym]{glossaries}
% \usepackage{acro}[=v2] % default is v3, and causing Overleaf timeout

%\makeglossaries

\DeclareAcronym{av}{   short=AV,  long=autonomous vehicle,  }
\DeclareAcronym{avs}{   short=AVs,  long=autonomous vehicles,  }

    % \acsetup{list-style=longtable}

    %pseudocode
    \usepackage{algorithm}
    \usepackage{algpseudocode}
    
    % various 
    \usepackage{booktabs}       % nicer tables
    \usepackage{multirow}
    \usepackage{todonotes}      % todo poznamky
    
    \usepackage{pmboxdraw}      % tree drawing
    \usepackage{verbatim}       % for verbatim environment and \verbatiminput
    \usepackage{sverb}          % for \verbinput
% GN setting end